{\KOMAoptions{open=any}%
\chapter*{Abstract}}

This thesis describes Aesop, a widely used proof automation tactic for the Lean~4 interactive theorem prover.
Its foundation is a tableau-style tree search similar to that of Rocq's and Isabelle's auto tactics, which users can customise by providing sets of proof rules that Aesop should apply.
On this foundation, we build three major novel features.

First, while similar tools use depth-first search, Aesop supports arbitrary search strategies and defaults to best-first search.
As a result, the common backtracking-based method for handling metavariables (also known as existential or schematic variables) during the search is no longer applicable.
We develop a novel method that preserves the independence of goals in the search tree even in the presence of metavariables while remaining, we conjecture, complete.

Second, we optimise the application of forward rules, which derive additional facts from local hypotheses.
Similar tools that support forward rules use a somewhat naive, stateless method to apply them, involving considerable redundant work.
We show how to efficiently cache some of the work that would otherwise be duplicated between similar goals.
This delivers speedups of over 30x in specific examples involving heavy use of forward rules, and a mean speedup of 1.24x on a large-scale benchmark where forward rules are more lightly used.
Like our metavariable handling, the new forward reasoning algorithm is not specific to dependent type theory and could be ported to systems based on different logics.

Third, we enable Aesop to generate largely idiomatic tactic proofs corresponding to the rules it applied.
This allows users to replace expensive Aesop searches with the generally much more efficient proofs Aesop found, or to use Aesop as a preprocessor that takes some obvious proof steps before handing control back to the user.
To make the generated proofs more idiomatic, we implement an optimiser that turns raw sequences of tactics into structured proofs and improves certain unidiomatic constructions.
